\section{Deep Dive: Why the Strongest External Tool Fails on This Benchmark}
\label{sec:tooldeepdive}
The head-to-head of the Results section deserves a closer look,
because the margin (84.13\% vs 38.46\% on the identical 624 official
test films) is large enough to demand an explanation rather than a
celebration. TorchXRayVision's DenseNet-121 is a serious, widely-cited
model; its failure mode here is diagnostic of the whole benchmark
landscape, and it is fully visible in the committed probability file
\texttt{test\_probs\_torchxrayvision.json}.
\subsection{The confusion structure}
At the 0.5 operating point the external model produces:
\begin{center}
\begin{tabular}{lcc}
\hline
 & Predicted NORMAL & Predicted PNEUMONIA \\
\hline
Actual NORMAL (234) & 234 & 0 \\
Actual PNEUMONIA (390) & 384 & 6 \\
\hline
\end{tabular}
\end{center}
Sensitivity is \textbf{1.5\%} (6/390) at 100\% specificity. The model
has effectively learned to say ``normal'' on this collection: its mean
predicted probability is 0.019 on actual pneumonia films and 0.0008 on
actual normals. A trivial all-negative classifier would score 37.5\%;
the 138M-parameter model scores 38.46\%.
\subsection{Why}
Three compounding causes, all structural rather than incidental:
\begin{enumerate}
\item \textbf{Prior shift.} TorchXRayVision's pneumonia head was trained
on adult collections (NIH ChestX-ray14, CheXpert, MIMIC) where pneumonia
prevalence is a few percent; this collection's test set is 62.5\%
pneumonia. A decision boundary calibrated to a rare disease cannot
survive a 20-fold prevalence increase without recalibration.
\item \textbf{Population shift.} The model is adult-trained; every film
in this benchmark is paediatric (Kermany et al.'s cohort). Paediatric
chest anatomy and pathology presentation differ enough that adult-tuned
features under-fire.
\item \textbf{Acquisition shift.} The films here are JPEG-compressed,
variably-sized exports, preprocessed by us to 128$\times$128 greyscale
--- far from the DICOM-native pipelines these models expect.
\end{enumerate}
None of this makes the external model bad; it makes the
\textbf{benchmark question} different from the deployment question. The
honest claim of this paper is narrow and verified: \emph{on this public
benchmark, under its official split, with no tuning of either side, our
compact models beat the strongest runnable published tool by 45.7
points, and the tool's published probabilities show why.} A deployment
claim would require the transfer study of the transfer study, which
we report honestly as a negative.
\subsection{What would change the verdict}
Threshold recalibration on the test prior (a single number, fit on the
test labels, hence not a fair comparison) lifts the external model
substantially; fine-tuning it on this training set would presumably
close most of the remaining gap at 1000$\times$ our parameter count.
We report the untuned comparison because it is the one a user
downloading both tools actually gets.
